% GENERATED FILE. Edit canonical lessons, then run npm run generate:books.
% canonical-source-hash: fnv1a64:dec1ab66f6d1bfb0
% canonical-lessons: MR-C420-dutavas, MR-C420-tarangan, MR-C420-pranisangrahalay, MR-C420-rahivasi, MR-C420-parisar

\chapter{Embassies and Zoos}
\label{ch:mr-embassies-and-zoos}
\hlchaptermodality{420}

\begin{chapteropening}
\textbf{By the end of this chapter:} \emph{I can talk about an embassy, a planetarium, a zoo, residents and the area around.}

Complete the last lesson of chapter 420: I can talk about an embassy, a planetarium, a zoo, residents and the area around.
\end{chapteropening}

\section[\texorpdfstring{\mr{दूतावास} (dūtāvās)}{दूतावास (dūtāvās)}]{\mr{दूतावास} (dūtāvās) — an embassy}

\label{lesson:MR-C420-dutavas}



\begin{quote}
Before the new one: say the Marathi for a narrow lane, then the Marathi for a tenement.
\end{quote}

\subsection*{You'll want to know: \mr{दूतावास}}
\textbf{\mr{दूतावास}} — \emph{dūtāvās} — \textquotedblleft{}an embassy\textquotedblright{}.

\begin{quote}
\mr{मीरा} \mr{व्हिसासाठी} \mr{दूतावासात} \mr{गेली}.
\end{quote}

\begin{quote}
\emph{mīrā vhisāsāṭhī dūtāvāsāt gelī.}
\end{quote}

\begin{quote}
Meera went to the embassy for her visa.
\end{quote}

Masculine.

\mr{दूत} is an envoy and \mr{आवास} a dwelling: the envoy’s house.

\begin{cousinweb}
\textbf{\mr{दूतावास}} (\emph{dūtāvās}) is a modern word built from Sanskrit \emph{dūta}, \textquotedblleft{}an envoy\textquotedblright{}, and \emph{āvāsa}, \textquotedblleft{}a dwelling\textquotedblright{} (Monier-Williams).
\end{cousinweb}

\begin{grammarlens}[title={What you've built}]
One more word for getting around.
\end{grammarlens}

\subsection*{Your turn}
\begin{itemize}
\raggedright
  \item \emph{Say it:} \emph{dūtāvās}
  \item \emph{Say it:} \emph{dūtāvās}, once more
  \item \emph{Say it:} \emph{mīrā vhisāsāṭhī dūtāvāsāt gelī.} — Meera went to the embassy for her visa.
  \item \emph{Recall:} say the Marathi for a narrow lane, then the Marathi for a tenement, then say \emph{dūtāvās} again
\end{itemize}

\begin{tcolorbox}[breakable,colback=teal!4,colframe=teal!35!black,title={Before you move on}]
What does \mr{दूतावास} mean? (\textquotedblleft{}An embassy\textquotedblright{}.) Say it once more.
\end{tcolorbox}
\section[\texorpdfstring{\mr{तारांगण} (tārāṅgaṇ)}{तारांगण (tārāṅgaṇ)}]{\mr{तारांगण} (tārāṅgaṇ) — a planetarium}

\label{lesson:MR-C420-tarangan}



\begin{quote}
Before the new one: say the Marathi for a tenement, then the Marathi for an embassy.
\end{quote}

\subsection*{You'll want to know: \mr{तारांगण}}
\textbf{\mr{तारांगण}} — \emph{tārāṅgaṇ} — \textquotedblleft{}a planetarium\textquotedblright{}.

\begin{quote}
\mr{शाळेची} \mr{सहल} \mr{तारांगणात} \mr{गेली}.
\end{quote}

\begin{quote}
\emph{śāḷecī sahal tārāṅgaṇāt gelī.}
\end{quote}

\begin{quote}
The school trip went to the planetarium.
\end{quote}

Neuter.

\mr{तारा}, a star, and \mr{अंगण}, a yard: a courtyard of stars.

\begin{cousinweb}
\textbf{\mr{तारांगण}} (\emph{tārāṅgaṇ}) is a modern word built from Sanskrit \emph{tārā}, \textquotedblleft{}a star\textquotedblright{}, and \emph{aṅgaṇa}, \textquotedblleft{}a courtyard\textquotedblright{} (Monier-Williams).
\end{cousinweb}

\begin{grammarlens}[title={What you've built}]
One more word for getting around.
\end{grammarlens}

\subsection*{Your turn}
\begin{itemize}
\raggedright
  \item \emph{Say it:} \emph{tārāṅgaṇ}
  \item \emph{Say it:} \emph{tārāṅgaṇ}, once more
  \item \emph{Say it:} \emph{śāḷecī sahal tārāṅgaṇāt gelī.} — The school trip went to the planetarium.
  \item \emph{Recall:} say the Marathi for a tenement, then the Marathi for an embassy, then say \emph{tārāṅgaṇ} again
\end{itemize}

\begin{tcolorbox}[breakable,colback=teal!4,colframe=teal!35!black,title={Before you move on}]
What does \mr{तारांगण} mean? (\textquotedblleft{}A planetarium\textquotedblright{}.) Say it once more.
\end{tcolorbox}
\section[\texorpdfstring{\mr{प्राणिसंग्रहालय} (prāṇisaṅgrahālay)}{प्राणिसंग्रहालय (prāṇisaṅgrahālay)}]{\mr{प्राणिसंग्रहालय} (prāṇisaṅgrahālay) — a zoo}

\label{lesson:MR-C420-pranisangrahalay}



\begin{quote}
Before the new one: say the Marathi for an embassy, then the Marathi for a planetarium.
\end{quote}

\subsection*{You'll want to know: \mr{प्राणिसंग्रहालय}}
\textbf{\mr{प्राणिसंग्रहालय}} — \emph{prāṇisaṅgrahālay} — \textquotedblleft{}a zoo\textquotedblright{}.

\begin{quote}
\mr{प्राणिसंग्रहालयात} \mr{आम्ही} \mr{वाघ} \mr{पाहिला}.
\end{quote}

\begin{quote}
\emph{prāṇisaṅgrahālayāt āmhī vāgh pāhilā.}
\end{quote}

\begin{quote}
We saw a tiger at the zoo.
\end{quote}

Neuter.

\mr{प्राणी}, an animal, shortens to \mr{प्राणि}- at the front of a compound.

\begin{cousinweb}
\textbf{\mr{प्राणिसंग्रहालय}} (\emph{prāṇisaṅgrahālay}) is a modern word built from Sanskrit \emph{prāṇin}, \textquotedblleft{}a living being\textquotedblright{}, \emph{saṁgraha}, \textquotedblleft{}a collection\textquotedblright{}, and \emph{ālaya}, \textquotedblleft{}a house\textquotedblright{} (Monier-Williams).
\end{cousinweb}

\begin{grammarlens}[title={What you've built}]
One more word for getting around.
\end{grammarlens}

\subsection*{Your turn}
\begin{itemize}
\raggedright
  \item \emph{Say it:} \emph{prāṇisaṅgrahālay}
  \item \emph{Say it:} \emph{prāṇisaṅgrahālay}, once more
  \item \emph{Say it:} \emph{prāṇisaṅgrahālayāt āmhī vāgh pāhilā.} — We saw a tiger at the zoo.
  \item \emph{Recall:} say the Marathi for an embassy, then the Marathi for a planetarium, then say \emph{prāṇisaṅgrahālay} again
\end{itemize}

\begin{tcolorbox}[breakable,colback=teal!4,colframe=teal!35!black,title={Before you move on}]
What does \mr{प्राणिसंग्रहालय} mean? (\textquotedblleft{}A zoo\textquotedblright{}.) Say it once more.
\end{tcolorbox}
\section[\texorpdfstring{\mr{रहिवासी} (rahivāsī)}{रहिवासी (rahivāsī)}]{\mr{रहिवासी} (rahivāsī) — a resident}

\label{lesson:MR-C420-rahivasi}



\begin{quote}
Before the new one: say the Marathi for a planetarium, then the Marathi for a zoo.
\end{quote}

\subsection*{You'll want to know: \mr{रहिवासी}}
\textbf{\mr{रहिवासी}} — \emph{rahivāsī} — \textquotedblleft{}a resident\textquotedblright{}.

\begin{quote}
\mr{या} \mr{इमारतीतले} \mr{रहिवासी} \mr{शांत} \mr{आहेत}.
\end{quote}

\begin{quote}
\emph{yā imāratītle rahivāsī śānt āhet.}
\end{quote}

\begin{quote}
The residents of this building are quiet.
\end{quote}

Masculine.

The plural is the same: \mr{एक} \mr{रहिवासी}, \mr{सगळे} \mr{रहिवासी}.

\begin{cousinweb}
\textbf{\mr{रहिवासी}} (\emph{rahivāsī}) is a Marathi word for someone who lives in a place; it ends like \textbf{\mr{प्रवासी}}, but how it was formed is not settled.
\end{cousinweb}

\begin{grammarlens}[title={What you've built}]
One more word for getting around.
\end{grammarlens}

\subsection*{Your turn}
\begin{itemize}
\raggedright
  \item \emph{Say it:} \emph{rahivāsī}
  \item \emph{Say it:} \emph{rahivāsī}, once more
  \item \emph{Say it:} \emph{yā imāratītle rahivāsī śānt āhet.} — The residents of this building are quiet.
  \item \emph{Recall:} say the Marathi for a planetarium, then the Marathi for a zoo, then say \emph{rahivāsī} again
\end{itemize}

\begin{tcolorbox}[breakable,colback=teal!4,colframe=teal!35!black,title={Before you move on}]
What does \mr{रहिवासी} mean? (\textquotedblleft{}A resident\textquotedblright{}.) Say it once more.
\end{tcolorbox}
\section[\texorpdfstring{\mr{परिसर} (parisar)}{परिसर (parisar)}]{\mr{परिसर} (parisar) — the surroundings; an area}

\label{lesson:MR-C420-parisar}



\begin{quote}
Before the new one: say the Marathi for a zoo, then the Marathi for a resident.
\end{quote}

\subsection*{You'll want to know: \mr{परिसर}}
\textbf{\mr{परिसर}} — \emph{parisar} — \textquotedblleft{}the surroundings; an area\textquotedblright{}.

\begin{quote}
\mr{मंदिराचा} \mr{परिसर} \mr{खूप} \mr{हिरवा} \mr{आहे}.
\end{quote}

\begin{quote}
\emph{mandirācā parisar khūp hirvā āhe.}
\end{quote}

\begin{quote}
The area around the temple is very green.
\end{quote}

Masculine.

\mr{शाळेचा} \mr{परिसर}, the school grounds and what is around them.

\begin{cousinweb}
\textbf{\mr{परिसर}} (\emph{parisar}) is a tatsama, taken from Sanskrit \emph{parisara}, \textquotedblleft{}neighbourhood, environs\textquotedblright{} (Monier-Williams).
\end{cousinweb}

\begin{grammarlens}[title={What you've built}]
One more word for getting around.
\end{grammarlens}

\subsection*{Your turn}
\begin{itemize}
\raggedright
  \item \emph{Say it:} \emph{parisar}
  \item \emph{Say it:} \emph{parisar}, once more
  \item \emph{Say it:} \emph{mandirācā parisar khūp hirvā āhe.} — The area around the temple is very green.
  \item \emph{Recall:} say the Marathi for a zoo, then the Marathi for a resident, then say \emph{parisar} again
\end{itemize}

\begin{tcolorbox}[breakable,colback=teal!4,colframe=teal!35!black,title={Before you move on}]
What does \mr{परिसर} mean? (\textquotedblleft{}The surroundings; an area\textquotedblright{}.) Say it once more.
\end{tcolorbox}
